% methodology.tex -- Cross-Enterprise Verification Paper (RU-V7)

\section{Methodology}

\subsection{Cross-Reference Circuit Design}

The cross-reference circuit proves the following statement in zero knowledge:
\emph{``There exist leaves in Enterprise $A$'s and Enterprise $B$'s sparse
Merkle trees, at positions determined by keys $k_A$ and $k_B$, whose values
satisfy an interaction predicate, and the Poseidon commitment of these values
matches the public interaction commitment.''}

\subsubsection{Public Inputs (3 field elements)}
\begin{enumerate}
\item $r_A$: Enterprise $A$'s state root (already public on L1 from prior
  batch submission)
\item $r_B$: Enterprise $B$'s state root (already public on L1)
\item $c$: Interaction commitment, $c = \text{Poseidon}(k_A, h_A, k_B, h_B)$
  where $h_X = \text{Poseidon}(k_X, v_X)$ is the leaf hash
\end{enumerate}

\subsubsection{Private Inputs (per interaction)}
For each enterprise $X \in \{A, B\}$:
\begin{itemize}
\item $k_X$: Leaf key (field element)
\item $v_X$: Leaf value (field element)
\item $\mathbf{s}_X = (s_{X,0}, \ldots, s_{X,31})$: Merkle sibling hashes
\item $\mathbf{b}_X = (b_{X,0}, \ldots, b_{X,31})$: Path direction bits
\end{itemize}

\subsubsection{Circuit Logic}

The circuit performs three operations:

\begin{enumerate}
\item \textbf{Merkle path verification} (per enterprise): Starting from leaf
  hash $h_X = \text{Poseidon}(k_X, v_X)$, iteratively hash with siblings
  according to path bits for 32 levels, and constrain the result to equal $r_X$.
  This is the \texttt{MerklePathVerifier} subcircuit from our prior
  work~\cite{ruv1smt}.

\item \textbf{Interaction commitment check}: Compute
  $c' = \text{Poseidon}(k_A, h_A, k_B, h_B)$ and constrain $c' = c$.

\item \textbf{Root binding}: The Merkle path verification inherently binds
  the proof to specific state roots $r_A$ and $r_B$, which are verified on L1
  as part of the public inputs.
\end{enumerate}

\subsubsection{Constraint Count}

\begin{table}[htbp]
\centering
\caption{Cross-Reference Circuit Constraint Breakdown}
\label{tab:constraints}
\begin{tabular}{@{}lr@{}}
\toprule
Component & Constraints \\
\midrule
Merkle path verification (side A, depth 32) & 34,254 \\
Merkle path verification (side B, depth 32) & 34,254 \\
Interaction predicate (Poseidon-4 + checks) & 360 \\
\midrule
\textbf{Total} & \textbf{68,868} \\
\bottomrule
\end{tabular}
\end{table}

The constraint count follows the formula $2 \times 1{,}038 \times (d+1) + 360$
where $d = 32$ is the tree depth and $1{,}038$ constraints per level is the
established rate from our Merkle path verifier~\cite{ruv2circuit}. This
circuit is approximately $2\times$ the size of a single-enterprise state
transition circuit (34,254 constraints), which is expected since it performs
two independent Merkle path verifications.

\subsection{Verification Approaches}

We evaluate three approaches for verifying cross-enterprise proofs on L1.

\subsubsection{Approach 1: Sequential Verification}

Each Groth16 proof is verified independently in separate transactions or
within a single transaction that calls the verifier contract multiple times.
For $N$ enterprises with $M$ cross-enterprise interactions:

\begin{equation}
G_\text{seq} = N \cdot G_\text{enterprise} + M \cdot G_\text{crossref}
\end{equation}

where $G_\text{enterprise} = 285{,}756$ gas (from RU-V3~\cite{ruv3gas}) and
$G_\text{crossref}$ is the Groth16 verification gas for 3 public inputs plus
cross-reference storage.

\subsubsection{Approach 2: Batched Pairing Verification}

All $N + M$ Groth16 proofs are batch-verified using a random linear
combination of pairing equations. Instead of $N + M$ independent pairing
checks (each requiring 4 pairings), a single combined pairing check is
performed with random coefficients~\cite{bellare1998batchrsa}. The dominant
savings come from replacing $(N + M - 1) \times 4$ pairings with scalar
multiplications:

\begin{equation}
G_\text{batch} = G_\text{pairing} + (N + M) \cdot G_\text{scalar} + G_\text{storage}
\end{equation}

This requires a hub coordinator that collects all proofs before submitting a
single batched verification transaction.

\subsubsection{Approach 3: Hub Aggregation}

An aggregation protocol (SnarkPack~\cite{gailly2022snarkpack} or Nebra
UPA~\cite{nebra2024upa}) compresses $N + M$ proofs into a single $O(\log(N+M))$
verification. We model this using the Nebra UPA gas formula:

\begin{multline}
G_\text{hub} = (N + M) \biggl(\frac{100K}{N+M} + 20K \\
  + \frac{350K}{N+M} + 7K + 25K\biggr) + G_\text{storage}
\end{multline}

Hub aggregation is most efficient at large proof counts (crossover at $\sim$32
proofs for SnarkPack) but incurs setup overhead at small counts.

\subsection{Privacy Analysis Methodology}

We analyze information leakage by enumerating all public signals and
determining what each reveals:

\begin{itemize}
\item $r_A, r_B$: Already public from individual enterprise submissions.
  \textbf{Additional leakage: 0 bits.}
\item $c = \text{Poseidon}(k_A, h_A, k_B, h_B)$: Under 128-bit preimage
  resistance, reveals nothing about the preimage. The act of submitting $c$
  reveals that an interaction exists. \textbf{Additional leakage: 1 bit}
  (interaction existence).
\item ZK proof $\pi$: Groth16 is zero-knowledge---the proof itself reveals
  nothing beyond the truth of the statement. \textbf{Additional leakage: 0 bits.}
\end{itemize}

Total information leakage per cross-enterprise interaction: \textbf{1 bit}
(interaction existence). This is the theoretical minimum: any protocol that
verifies cross-enterprise interactions on a public ledger must submit
\emph{something}, which inherently reveals that a verification occurred.
